\documentclass{article}

\usepackage{fullpage}
\usepackage{url}
\usepackage{graphicx}
\usepackage{amsmath}
\usepackage{physics}
\usepackage{mathtools}
\usepackage{bbold}

\usepackage{lmodern}
\usepackage[T1]{fontenc}
\usepackage[fontsize=13.5pt]{fontsize}

\DeclareMathOperator{\erfi}{erfi}
\newcommand{\R}{\ensuremath{\mathbb{R}}}
\newcommand{\C}{\ensuremath{\mathbb{C}}}

\title{Complex Systems Excercise 1}
\author{Idan Stark}
\date{\today}

\begin{document}

\maketitle

\section{Dimensional analysis - Flow in porous media}
Soap films are fluid membranes, whose shape in space is dominated by surface tension. As such, an unconstrained soap film (not enclosing any volume) will form a minimal surface.
\newline
Let us consider two identical rings of radius $R$ dipped together in soap-water, and then drawn apart from each other. Assuming they make an open surface, and assuming the symmetry of the problem is maintained, let us look for a solution in the form of a surface of revolution:
\begin{equation}
    \vec{r}(z, \nu) = \phi(z)\cos(\nu)\hat{x} + \phi(z)\sin(\nu)\hat{y} + z\hat{z}
\end{equation}
where $\nu \in [0, 2\pi]$ and $\phi(z)$ is the profile function.
\subsection{Tangent vectors}
The tangent vectors of the surface are given by:
\begin{equation}
    \boxed{
\begin{gathered}
        \vec{r}_z = \partial_z \vec{r} = \phi'(z)\cos(\nu)\hat{x} + \phi'(z)\sin(\nu)\hat{y} + \hat{z}
        \\
        \vec{r}_\nu = \partial_\nu \vec{r} = -\phi(z)\sin(\nu)\hat{x} + \phi(z)\cos(\nu)\hat{y}
\end{gathered}
    }
\end{equation}
\subsection{The metric tensor}
This gives:
\begin{equation}
\begin{gathered}
    a_{zz} = \vec{r}_z\cdot\vec{r}_z = \phi'^2(z)\cos^2(\nu) + \phi'^2(z)\sin^2(\nu) + 1^2 = 1 + \phi'^2(z)
    \\
    a_{\nu z} = a_{z \nu} = \vec{r}_\nu\cdot\vec{r}_z = -\phi(z)\phi'(z)\cos(\nu)\sin(\nu) + \phi(z)\phi'(z)\cos(\nu)\sin(\nu) + 0 = 0
    \\
    a_{\nu\nu} = \vec{r}_\nu\cdot\vec{r}_\nu = \phi^2(z)\sin^2(\nu) + \phi^2(z)\cos^2(\nu) = \phi^2(z)
\end{gathered}
\end{equation}
So we get:
\begin{equation}
    \boxed{a_{\alpha\beta} = \begin{pmatrix}
        1 + \phi'^2 & \\ & \phi^2
    \end{pmatrix}}
\end{equation}
\subsection{The bending tensor and shape tensor}
First, let us find the normal vector:
\begin{equation}
    \hat{N} = \frac{\vec{r}_z \times \vec{r}_\nu}{\sqrt{a}} = \frac{-\phi(z)\cos(\nu)\hat{x} - \phi(z)\sin(\nu)\hat{y} + \phi(z)\phi'(z)\hat{z}}{\sqrt{(1 + \phi'^2)\phi^2}}
\end{equation}
Also:
\begin{equation}
\begin{gathered}
    \partial_{z}\partial_{z} \vec{r} = \phi''(z)\cos(\nu)\hat{x} + \phi''(z)\sin(\nu)\hat{y}
    \\
    \partial_{z}\partial_{\nu} \vec{r} = \partial_{\nu}\partial_{z} \vec{r} = -\phi'(z)\sin(\nu)\hat{x} + \phi'(z)\cos(\nu)\hat{y}
    \\
    \partial_{\nu}\partial_{\nu} \vec{r} = -\phi(z)\cos(\nu)\hat{x} - \phi(z)\sin(\nu)\hat{y}
\end{gathered}
\end{equation}
Which gives:
\begin{equation}
    \boxed{
        b_{\alpha\beta} = \partial_\alpha\partial_\beta \vec{r}\cdot\hat{N} =
        \frac{1}{\sqrt{1 + \phi'^2}}
        \begin{pmatrix}
            -\phi'' & \\ & \phi
        \end{pmatrix}
    }
\end{equation}
And the shape tensor is:
\begin{equation}
    s^\alpha_\beta =
    \frac{1}{\sqrt{1 + \phi'^2}}
    \begin{pmatrix}
        \frac{1}{1 + \phi'^2} & \\ & \frac{1}{\phi^2}
    \end{pmatrix}
    \begin{pmatrix}
        -\phi'' & \\ & \phi
    \end{pmatrix} =
    \boxed{\frac{1}{\sqrt{1 + \phi'^2}}
    \begin{pmatrix}
        -\frac{\phi''}{1 + \phi'^2} & \\ & \frac{1}{\phi}
    \end{pmatrix}}
\end{equation}
\subsection{Minimal surface}
For a minimal (open) surface, we reqiure
\begin{equation}
    H = 0 \rightarrow \Tr{s} = 0
\end{equation}
This gives the equation:
\begin{equation}
    \boxed{-\frac{\phi''}{1 + \phi'^2} + \frac{1}{\phi} = 0}
\end{equation}
\subsection{Solving the ODE}
Multiplying and organizing gives:
\begin{equation}
    \phi\phi'' - \phi'^2 - 1 = 0
\end{equation}
Using:
\begin{equation}
    \phi'' = \frac{d\phi'}{dx} = \frac{d\phi'}{d\phi}\frac{d\phi}{dx} = \phi'\frac{d\phi'}{d\phi} = \frac{1}{2}\frac{d}{d\phi}\phi'^2
\end{equation}
We can define the derivative as a function of $\phi$ to get $\sigma(\phi) = \phi'^2(\phi)$. The ODE the reduces to a linear ODE:
\begin{equation}
    \phi\frac{d\sigma}{d\phi} - 2\sigma - 2 = 0
\end{equation}
Seperating:
\begin{equation}
    \frac{d\sigma}{2\sigma + 2} = \frac{d\phi}{\phi}
\end{equation}
Integrating:
\begin{equation}
    \frac{1}{2}\ln(\sigma + 1) = \ln\phi + C
\end{equation}
And organizing:
\begin{equation}
    \sigma = A \phi^2 - 1
\end{equation}
Then pluggin back the definition:
\begin{equation}
    \phi' = \sqrt{A\phi^2 - 1}
\end{equation}
This is simpler equation, but to make it even easier let us find $A$ using the boundary conditions on $z = 0$. This will give $\phi' = 0, \phi = r_{min}$ and thus:
\begin{equation}
    Ar_{min}^2 - 1 = 0 \rightarrow A = r_{min}^{-2}
\end{equation}
So, we can define $\varphi = \phi/r_min$ to get:
\begin{equation}
    r_{min}\varphi' = \sqrt{\varphi^2 - 1}
\end{equation}
Using seperation of variables again:
\begin{equation}
    r_{min}\frac{d\varphi}{\sqrt{\varphi^2 - 1}} = dz
\end{equation}
And integrating
\begin{equation}
    r_{min} \cosh^{-1}\varphi = z + B
\end{equation}
And finally solving for $\phi$ gives:
\begin{equation}
    \phi = r_{min}\cosh(\frac{z + B}{r_{min}})
\end{equation}
Pluggin $\phi(0) = r_{min}$ gives $B = 0$, and so we get:
\begin{equation}
    \boxed{\phi = r_{min}\cosh(\frac{z}{r_{min}})}
\end{equation}
\subsection{Resutling metricand shape tensors}
We can blue this result into the metric and shape tensors, and use $1 - \sinh^2(x) = \cosh^2(x)$ to get:
\begin{equation}
\begin{gathered}
    a_{\alpha\beta} = 
    \begin{pmatrix}
        1 + \sinh^2(z/r_{min}) & \\ & r_{min}^2\cosh^2(z/r_{min})
    \end{pmatrix} = \\ =
    \cosh^2(z/r_{min})
    \begin{pmatrix}
        1 & \\ & r_{min}^2
    \end{pmatrix}
\end{gathered}
\end{equation}
\begin{equation}
\begin{gathered}
    s^\alpha_\beta =
    \frac{1}{\sqrt{1 + \sinh^2(z/r_{min})}}
    \begin{pmatrix}
        -\frac{\cosh(z/r_{min})}{r_{min}\left(1 + \sinh^2(z/r_{min})\right)} & \\ & \frac{1}{r_{min}\cosh(z/r_{min})}
    \end{pmatrix} = \\ =
    \frac{1}{r_{min}\cosh^2(z/r_{min})}
    \begin{pmatrix}
        -1 & \\ & 1
    \end{pmatrix}
\end{gathered}
\end{equation}
To summarize:
\begin{equation}
    \boxed{a_{\alpha\beta} = \cosh^2(z/r_{min})
    \begin{pmatrix}
        1 & \\ & r_{min}^2
    \end{pmatrix}
    \qquad
    s^\alpha_\beta =
    \frac{1}{r_{min}\cosh^2(z/r_{min})}
    \begin{pmatrix}
        -1 & \\ & 1
    \end{pmatrix}
    }
\end{equation}
\subsection{Given rings}
If the two rings have a given radius $R$, and are placed at $z = \pm z_0$, we can plug this into the $\phi(z)$ to get:
\begin{equation}
    \boxed{r_{min}\cosh(z_0/r_{min}) = R}
\end{equation}
\subsection{Explaining the video}
In the video we see the shape that we calculated. The rings start with $z_0 \ll R$, and then as they grow appart and $z_0$ grows, to keep $R$ constant, $r_{min}$ must shrink.
At a certain point, the surface width becomes too small to sustain, then the bubble pops.

\section{Buckling}
\subsection{Assumption}
In this section we use the dimensionally reduced description we saw in class. It includes the following assumptions:
\begin{itemize}
    \item Plane stress assumption: $W$ is small and as such $\sigma^{i3} = 0$. This is basically "there is no place for stress in the $w$ dimension".
    \item Thin body assumption: $t$ is small and as such we can ignore deformations in it across the body. This is done through a similar assumption to the first one: $\sigma^{\alpha 2} = 0$, and the assumption that by going along the tangent vector to the midline we can get to any point in the configuration:
    \begin{equation}
        \vec{r}(u, v) = \vec{\gamma}(u) + v\hat{n}(u)
    \end{equation}
    where $\vec{\gamma}$ is the midline curve and $\hat{n}(u)$ is the normal to the curve.
\end{itemize}
\subsection{Pure bending mode}
In the pure bending mode, there is no stretching, and thus:
\begin{equation}
    \epsilon^{11} = 0
\end{equation}
\subsubsection{Geometric relation}
The ruler will bend to create an arc, classified by a certain openning angle $\theta$ and radius $R$. Since there is no stretching, the length of the arc is the same as before:
\begin{equation}
    L = R\theta
\end{equation}
while the length of the cord is:
\begin{equation}
    L - \Delta = 2R \sin(\frac{\theta}{2})
\end{equation}
Dividing the two gives:
\begin{equation}
    \boxed{1 - \frac{\Delta}{L} = \frac{\sin(\theta/2)}{\theta/2}}
\end{equation}
\subsubsection{Small displacement approximation}
Now, let us take $\Delta \ll L$. This will result in the LHS in the above relation to be close to 1, and thus the right hand side will be close to 1. Since $\frac{\sin x}{x} \approx 1 \rightarrow x \approx 0$, we can expand the angle $\theta/2$ around zero to get:
\begin{equation}
    1 - \frac{\Delta}{L} \approx \frac{\theta/2 - \frac{(\theta/2)^3}{3!}}{\theta/2} = 1 - \frac{\theta^2}{24}
\end{equation}
\begin{equation}
    \boxed{\frac{\Delta}{L} = \frac{\theta^2}{24}}
\end{equation}
We can thus find the curvature:
\begin{equation}
    \kappa = R^{-1} = \frac{\theta}{L} = \sqrt{\frac{24\Delta}{L^3}}
\end{equation}
Which is constant. We can then write the bending energy:
\begin{equation}
    E_{B} = \frac{YWt^3}{24}\int_0^L \kappa^2 du = \boxed{\frac{YWt^3\Delta}{L^2}}
\end{equation}
\subsection{Condition for pure bending deformation}
we require that the pure bending mode will require less energy then the stress mode, which, given the constant deformation $\Delta$, has energy:
\begin{equation}
    E_{S} = \frac{YWt}{2L}\Delta^2
\end{equation}
Comparing the two gives:
\begin{equation}
    E_{B} < E_{S}
\end{equation}
\begin{equation}
    \frac{YWt^3\Delta}{L^2} < \frac{YWt}{2L}\Delta^2
\end{equation}
\begin{equation}
    \boxed{\Delta > \frac{2t^2}{L}}
\end{equation}
\end{document}